\subsection{What International AI Collaboration Can and Cannot Fix}
\label{sec:8.7.3}
AI collaboration reaches furthest where the underlying problem is inherently cross-border and no single country's data is sufficient on its own. Climate is the clearest case: weather and climate systems do not respect national boundaries, and a model trained only on one region's data predicts poorly everywhere else. The institutional machinery here is real and worth describing precisely, because descriptions of it are usually inflated. The Climate, Land, Energy and Water Systems framework — coordinated by UN bodies together with the International Atomic Energy Agency, with a capacity-building program that has reached participants in more than twenty countries — is an integrated systems-optimization framework, and the research literature around it treats machine learning as a promising addition still being built in, not a feature the framework already has \autocite{alexander2025clews}. Where AI collaboration does the specific climate work this section is about, it currently looks more like national weather services training and benchmarking forecasting models against shared data than like a single named research consortium.

Disease surveillance is a similar case for the same underlying reason: a pathogen crossing from an animal population into a human one, or from one country into the next, does not stay inside the jurisdiction that first detects it. The Global Virome Project, proposed as a successor to USAID's PREDICT program, set out to catalog the viral families with the highest zoonotic spillover potential across dozens of countries before the next pandemic instead of after it \autocite{carroll2018virome}. The AI in that picture sits downstream of the cataloging. Sequence data of the kind such a project generates is what spillover-risk prediction models get trained on, which makes the international data-sharing arrangement itself, and not any single algorithm, the part actually worth calling collaboration.

Poverty and economic development are where this kind of cross-border AI collaboration is least developed. The general claims are easy to make — AI-driven analysis revealing patterns of inequality that guide where aid and investment go, better agricultural forecasting, financial services reaching places that have none. A specific, verifiable international partnership doing sustained work at that scale is harder to find than those claims suggest. What exists at the institutional level is the digital-divide problem section~\ref{sec:8.7.6} covers in detail: uneven data, uneven compute, uneven regulatory capacity, and a recurring set of proposed remedies, capacity building and technology transfer among them, that has not yet closed the gap.

Conflict prevention is a fourth domain where the underlying data problem — dispersed, multilingual, and produced by people with no interest in cooperating with researchers — makes international collaboration closer to a precondition than an option. PeaceTech Lab, an independent nonprofit that grew out of a US Institute of Peace project, is a working example: its hate-speech lexicon project has combined machine learning with local-language workshops and interviews to build region-specific glossaries of inflammatory language in at least eleven countries, including Kenya, Nigeria, Iraq, and Yemen — ground-level data collaboration that international peacebuilding tools depend on and rarely get \autocite{peacetechlab2019monitoring}.

None of these four cases is evidence that AI collaboration works the way the treaty layer aspires to. Each one ran on the researchers and institutions involved choosing to participate, project by project, with no enforcement mechanism binding the next one to happen and no guarantee that a government hostile to the work could not simply keep its own researchers out of it. Section~\ref{sec:6.4.1} catalogs what a government willing to point AI at its own population does with tools built for exactly these humanitarian purposes elsewhere. A government can fund pandemic-surveillance research of the kind described above and, separately, build the biometric-surveillance systems that section documents, and nothing about participating in the first constrains the second.
